% !TeX program = xelatex
% =====================================================================
% 陕西师范大学本科毕业论文（设计）开题报告 —— 通用模板
\newif\ifoverleaf
\overleaftrue
\ifoverleaf
  \documentclass[12pt,a4paper,fontset=fandol]{ctexart}
\else
  \documentclass[12pt,a4paper,fontset=windows]{ctexart}
\fi
\usepackage{amsmath,amssymb}
\usepackage{cite}
\usepackage{array}
\usepackage{multirow}
\usepackage{graphicx}
\usepackage{xcolor}
\usepackage{geometry}
\geometry{top=3.6cm,bottom=2.5cm,left=3.0cm,right=3.0cm}
\usepackage{fancyhdr}
\fancyhf{}
\fancyhead[C]{\zihao{5} 陕西师范大学本科毕业论文（设计）开题报告}
\fancyfoot[C]{\thepage}
\renewcommand{\headrulewidth}{0.4pt}
\setlength{\headheight}{15pt}
\setlength{\headsep}{18pt}
% ---- 正文一级标题（一、二、…）:四号加粗 ----
\newcommand{\kwsec}[1]{\par\vspace{1.1em}\noindent{\bfseries\zihao{4}#1}\par\vspace{0.35em}}
\begin{document}
% ============================ 封 面 ============================
\pagestyle{empty}
\setcounter{page}{0}  %%% 关键：封面将page计数器置0，封面不占页码
\begin{center}
  \vspace*{2.48cm}
  \includegraphics[width=11.56cm]{shanxishida_logo.jpg}
\end{center}
\vspace{0.25cm}
\begin{center}
  {\kaishu\zihao{1} 本科生毕业论文（设计）开题报告}
\end{center}
\vspace{1.66cm}
% ---------------- 题目（三号，楷体；可占一行或两行） ----------------
% 第一行写不下时换行续写；只需一行时，删除下面「题目第二行」那一行即可。
{\zihao{3}\noindent\hspace*{56pt}\heiti 题目\kaishu ：\hspace{0.5em}题目第一行（写不下的内容放第二行）\par}
\vspace{7.6pt}
{\zihao{3}\noindent\kaishu 题目第二行（若只有一行，删去本行）\par}
\vspace{2.26cm}
% ---------------- 作者信息（四号黑体，标签两端分散） ----------------
\begin{center}
	\renewcommand{\arraystretch}{1.64}
	\setlength{\tabcolsep}{6pt}
	{\zihao{4}\heiti
		\begin{tabular}{
				p{6.2cm}
				>{\centering\arraybackslash}p{5.8cm}
			}
			\makebox[\linewidth][s]{作 者 单 位} & 陕西师范大学 \\
			\makebox[\linewidth][s]{作 者 姓 名} & 张三 \\
			\makebox[\linewidth][s]{专 业 班 级} & 如“恒元2301班” \\
			\makebox[\linewidth][s]{作 者 学 号} & 学号 \\
			\makebox[\linewidth][s]{指 导 教 师（职 称）} & 如“XXX 教授”，无则留空 \\
	\end{tabular}}
\end{center}
\vspace{0.38cm}
% ---------------- 日期（四号黑体） ----------------
\begin{center}
  {\zihao{4}\heiti \  年\  月}
\end{center}
\clearpage   % 翻页，进入正文页，此时page计数器自动+1变为1
\pagestyle{fancy}
\linespread{1.5}\selectfont
% ============================ 一、选题背景与意义 ============================
\kwsec{一、选题背景与意义}
% 说明：本课题的学科背景、国内外研究现状与不足之处、拟研究问题、
%       理论意义与应用价值。可分 2—3 段写，引用处用 \cite{key} 并在文末列出。
%【正文填写处】
% ============================ 二、研究的主要目标和内容 ============================
\kwsec{二、研究的主要目标和内容}
% 说明：先写研究目标（总目标 + 分点目标），再分条写主要研究内容（（1）（2）（3）…）。
%【正文填写处】
% ============================ 三、拟采用的研究方法、步骤 ============================
\kwsec{三、拟采用的研究方法、步骤}
% 说明：研究方法（解析/数值/实验手段），可行性论证（方法是否成熟、数据/设备是否可得），
%       以及分阶段研究步骤（1）—（6）。
%【正文填写处】
% ============================ 四、课题研究进度安排 ============================
\kwsec{四、课题研究进度安排}
% 说明：下表为示例，请按自己实际年限与周次改列。多/少一个学期时，
%       复制/删除对应的 \multirow{4}{...} 块即可（并同步修改行数与 \cline 个数）。
\begin{center}
  {\bfseries\zihao{5} 表1\quad 课题研究进度安排表}
\end{center}
\renewcommand{\arraystretch}{1.5}
\setlength{\tabcolsep}{3pt}
\begin{center}
	\begin{tabular}{
			|>{\centering\arraybackslash}p{3.6cm}
			|>{\centering\arraybackslash}p{3.0cm}
			|>{\centering\arraybackslash}p{7.5cm}|
		}
		\hline
		\bfseries 学期 & \bfseries 周次 & \bfseries 工作任务 \\ \hline
		\multirow{4}{*}{\begin{tabular}{@{}c@{}}20XX—20XX学年\\第一学期\end{tabular}}
		& 第1—3周 & 明确选题，文献调研与开题准备 \\ \cline{2-3}
		& 第4—8周 & 撰写并提交开题报告 \\ \cline{2-3}
		& 第9—13周 & 核心研究内容的建模与实现 \\ \cline{2-3}
		& 第14—17周 & 阶段性计算/实验，中期自查 \\ \hline
		\multirow{4}{*}{\begin{tabular}{@{}c@{}}20XX—20XX学年\\第二学期\end{tabular}}
		& 第1—5周 & 第二学期第一阶段工作 \\ \cline{2-3}
		& 第6—9周 & 第二学期第二阶段工作 \\ \cline{2-3}
		& 第10—14周 & 撰写毕业论文初稿 \\ \cline{2-3}
		& 第15—17周 & 论文修改、查重与定稿，准备答辩 \\ \hline
	\end{tabular}
\end{center}
% ============================ 五、主要参考文献 ============================
\kwsec{五、主要参考文献}
% 说明：GB/T 7714 格式。中文在前、英文在后各按字母/拼音排序或按正文引用顺序。
%       示例两条如下，请按实际文献增删 \bibitem 条目并保留 \cite{key} 对应。
%       编号用 \bibitem 上标数字时，把 {12}（最大编号）改成实际条数。
%       可自行调整文献部分，适配不同的文献格式，这里展示的为最基础的文献格式。
\begingroup
\renewcommand{\section}[2]{}%   吞掉 thebibliography 自动生成的标题，避免与上文重复
\begin{thebibliography}{12}
\bibitem{exampleZh} 作者名. 中文文献题名[J]. 期刊名, 年份, 卷(期): 起止页码.
\bibitem{exampleEn} Author A B, Author C D. Title of the paper[J]. Journal Name, Year, Volume(Issue): Pages.
\end{thebibliography}
\endgroup
% =====================================================================
% 附：开题报告评审表
%   对应官方 Word 版末页《陕西师范大学本科生毕业论文（设计）开题报告评审表》，
%   排在“五、主要参考文献”之后、另起一页。可打印后手写，也可电子填写。
% 说明：
%   ① 本表已按官方 Word 版还原为真实 6 列结构（列宽对应 Word 的
%      998/1417/992/2410/1418/1770 twips），含单元格合并（\multicolumn）
%      与独立填值格。电子填写时把内容直接填进空白格即可。
%   ② 两处 \vspace{…cm} 控制“评语 / 委员会意见”手写留白的高度，按需增删。
%   ③ 单元格正文五号、加粗标题；红字提示用楷体红色。标题三号黑体；
%      若学校要求“华文中宋”加粗，把 \heiti 改成 \songti\bfseries。
\clearpage
\begin{center}
  {\songti\bfseries\zihao{3} 陕西师范大学本科生毕业论文（设计）开题报告评审表}
\end{center}
\vspace{0.3cm}

\begingroup
\setlength{\parindent}{0pt}
\renewcommand{\arraystretch}{1}
\renewcommand{\baselinestretch}{1}\selectfont
\setlength{\tabcolsep}{0pt}
% 列宽按 Word 的 998/1417/992/2410/1418/1770 twips 等比放大到总宽15.9cm，
% 并比正文(15.0cm)稍宽，用 \makebox 居中、左右各轻微出血 0.45cm，与 Word 版一致。
\noindent\makebox[\linewidth][c]{%
\begin{tabular}{
	| >{\centering\arraybackslash}m{1.75cm}
	| >{\centering\arraybackslash}m{2.49cm}
	| >{\centering\arraybackslash}m{1.75cm}
	| >{\centering\arraybackslash}m{4.20cm}
	| >{\centering\arraybackslash}m{2.49cm}
	| >{\centering\arraybackslash}m{3.12cm} |}
  \hline
  {\zihao{5}\bfseries 姓名} & \rule{0pt}{1.15cm} & {\zihao{5}\bfseries 学号} & & {\zihao{5}\bfseries 指导教师} & \\ \hline
  \multicolumn{2}{| >{\centering\arraybackslash}m{4.24cm}|}{\zihao{5}\bfseries 学院、专业}
  & \multicolumn{4}{ >{\centering\arraybackslash}m{11.56cm}|}{\rule{0pt}{1.3cm}} \\ \hline
  \multicolumn{6}{| p{15.80cm}|}{%
      \vspace{0.20cm}
      {\centering\zihao{5}\bfseries 指导教师评语\par}
      \vspace{0.30cm}
      {\zihao{5}
        \leftskip=0.35cm
        1．学生前期表现情况。\par
        2．是否具备开始设计（论文）条件？是否同意开始设计（论文）？\par
        3．不足及建议。\par}%
      \rule{0pt}{1.1cm}} \\ \hline
  \multicolumn{6}{| p{15.80cm}|}{%
      {\centering\kaishu\zihao{5}\color{red}（用黑色签字笔手写或小4号宋体字编辑，签名必须手写。可加页，A4纸双面打印）\par}
      \vspace{5.15cm}
      {\zihao{5}\hfill 指导教师（签名）：\hspace*{1.5cm}\par}
      \vspace{0.55cm}
      {\zihao{5}\hfill 年\qquad\qquad\qquad 月\qquad\qquad\qquad 日\hspace*{0.5cm}\par}} \\ \hline
  \multicolumn{6}{| >{\centering\arraybackslash}p{15.80cm}|}{%
      \vspace{0.20cm}
      {\zihao{5}\bfseries 学院答辩委员会意见\par}\rule{0pt}{0.55cm}} \\ \hline
  \multicolumn{6}{| p{15.80cm}|}{%
      \vspace{3.4cm}
      {\zihao{5}\hfill 负责人（签名）：\hspace*{1.5cm}\par}
      \vspace{0.55cm}
      {\zihao{5}\hfill 年\qquad\qquad\qquad 月\qquad\qquad\qquad 日\hspace*{0.5cm}\par}} \\ \hline
  \end{tabular}}
\endgroup
\end{document}
